\section{The envelope of admissible completions}
\label{sec:vsenvelope}

The ceiling constrains slopes at infinity.  A desk's question is more
concrete: the quotes end at some strike --- what prices may an
admissible model print at the strikes just past it?  This section answers
with an envelope: two curves, computable from the last two quotes,
between which every admissible continuation must run.

Fix the setting on the call side (the put side mirrors it).  A slice
has been fitted through the quotes; call everything up to the last
quoted strike \emph{the belly} and regard it as settled.  Bars mark
quantities frozen at the last quote: $\bar k$ is the last quoted
log-strike, $\bar y=e^{\bar k}$ the strike itself, and
$\bar c=c(\bar y)$ the normalized call there.  Write $y_1$ for the
second-to-last quoted strike; the \emph{last secant slope} is the
price change per unit of strike across the final quoted interval,
\[
\bar s=\frac{c(\bar y)-c(y_1)}{\bar y-y_1}\;<\;0 .
\]  A
\emph{completion} is any way of continuing the call curve to
strikes beyond $\bar y$ so that the whole curve still prices a law.
Chapter~2 established what that requires
(\cref{sec:intro,eq:BL}): as a function of strike, the call curve of
a law is nonincreasing, and it is convex --- its second derivative
\emph{is} the density, so convexity is positivity of probability.
And by \eqref{eq:vscalldies} it must eventually reach zero.

Those three properties already imply the envelope.

\begin{proposition}[The admissible cone]
\label{prop:vscone}
Every admissible completion satisfies, for all $y\ge\bar y$,
\begin{equation}
\max\!\big(0,\;\bar c+\bar s\,(y-\bar y)\big)
\;\le\;c(y)\;\le\;\bar c .
\label{eq:vscone}
\end{equation}
The lower edge is itself an admissible completion; the upper edge is
approached by admissible completions but attained by none.
\end{proposition}

\begin{proof}
Write the completion through its slope: $c(y)=\bar c+\int_{\bar
y}^{y}c'(t)\,\dd t$.  Convexity of the whole curve says $c'$ is
nondecreasing; in particular, for a convex function a chord's slope
never exceeds the slope just past its right endpoint, so applying
this to the chord over $[y_1,\bar y]$ gives $c'(\bar y^{+})\ge\bar
s$, hence $c'(t)\ge\bar s$ for all $t\ge\bar y$.  Monotonicity says
$c'\le0$.  Integrating the two slope bounds gives
$\bar c+\bar s(y-\bar y)\le c(y)\le\bar c$, and prices are
nonnegative, which is \eqref{eq:vscone}.

For the lower edge, let the slope stay exactly at $\bar s$ until the
line reaches zero, at the strike
\begin{equation}
y_0=\bar y-\frac{\bar c}{\bar s},
\label{eq:vsyzero}
\end{equation}
and stay at zero after.  Slopes step from $\bar s$ up to $0$, so the
curve is convex and nonincreasing, and it reaches zero: admissible.
Its kink at $y_0$ is a point mass --- the law parks all remaining
upside probability at the single strike $y_0$ and owns nothing
beyond.  For the upper edge, a concrete family does the work.  For
small $\alpha>0$ consider the completion
$c(y)=\bar c\,(y/\bar y)^{-\alpha}$: it is decreasing and convex, it
tends to zero, and it joins the belly with a convex kink at $\bar y$
(its opening slope $-\alpha\bar c/\bar y$ is closer to zero than
$\bar s$, and slopes may jump upward at a kink).  Being decreasing,
convex and vanishing at infinity, it is the call curve of a genuine
law (\cref{sec:intro}), one that pushes its remaining mass to ever
higher strikes as $\alpha$ shrinks.  As $\alpha\to0$ these curves
approach the constant $\bar c$ at every strike --- but the constant
itself never reaches zero, so it is not admissible.
\end{proof}

Pause on what the proposition says.  At any strike beyond the last
quote, the admissible call prices fill an interval: from a linear
descent that hits zero --- at $y_0$, which \eqref{eq:vsyzero}
computes from two quotes --- up to (but excluding) the last quoted
price itself.  On the running node the numbers are stark: the last
call-side quote sits at $\bar k=\MacWingEnvKbar$, worth
$\bar c=\MacWingEnvCbar$ of the forward, and the lower edge reaches
zero at $k=\MacWingEnvKzero$.  The data permits the smile to
\emph{end} --- price zero, no law beyond --- less than a tenth of a
log-unit past the last quote.  It also permits a wing that is still
worth nearly the full $\bar c$ a log-unit out.  Everything between is
allowed.

\begin{figure}[htbp]
\centering
\paperfig{\textwidth}{fig_vs_envelope.pdf}
\caption{The envelope of admissible completions (running node, call
side).  (a)~Price space: the fitted belly (black) through the last
quote mids (orange), and the cone \eqref{eq:vscone} beyond the last
quote: the dashed monotone cap at $\bar c$, the convex floor
descending at the secant slope $\bar s$, reaching zero at
$y_0$ (dot).  (b)~The same cone mapped through the Black chart: the
shaded fan $[w^{-},w^{+}]$ of admissible total-variance completions,
with the three families' fitted wings threading it.}
\label{fig:vsenvelope}
\end{figure}

\Cref{fig:vsenvelope}(a) draws the cone in price space on the running
node.  The black curve is the fitted belly, ending at the last quote;
prices here are a few $10^{-4}$ of the forward, so the axis is in
those units.  From the endpoint two edges open: the dashed horizontal
--- the \emph{monotone cap}, prices may not rise --- and the
descending line --- the \emph{convex floor}: prices may never fall
faster than the last secant slope, so the straight continuation at
$\bar s$, dying into the axis at $y_0$ (dot), is the steepest descent
any completion is allowed.  The shaded region between them is every
call price any law consistent with the belly can print.

Panel~(b) shows the same cone where this book draws smiles: in total
variance.  The Black function is increasing in $w$, so the price
interval \eqref{eq:vscone} maps strike by strike onto an interval
$[w^{-}(k),\,w^{+}(k)]$, the \emph{fan}: $w^{-}$ is the implied
variance of the descending line (zero once the line dies), $w^{+}$
that of the constant $\bar c$.  The fan is vast.  By $k=1$ the
admissible total variance runs from $0$ to several times the ATM
level; the three families' wings --- the same three fits of
\cref{fig:vspins} --- thread it in a narrow bundle near the bottom.
Numerically, each family's wing clears the floor and cap of its
\emph{own} cone --- the cone rebuilt from that family's fitted call
values at the same two strikes, since \cref{prop:vscone} constrains a
completion relative to its own belly (\cref{app:vsprotocol}) --- with
room to spare: worst clearances
\MacWingEnvClearLqdLo/\MacWingEnvClearSviLo/\MacWingEnvClearMcsLo\
below and
\MacWingEnvClearLqdHi/\MacWingEnvClearSviHi/\MacWingEnvClearMcsHi\
above, in total variance.  That is \cref{prop:vscone} confirmed on
real fits.
The bundle's narrowness against the fan's width restates
\cref{sec:vslives}: the models' structural choices are far more
selective than the data's constraints.

The fan's top edge has a closed form, and deriving it closes the loop
with \cref{sec:vsceiling}.  The upper edge holds the call at the
constant $\bar c$, so $w^{+}(k)$ solves
$\Black(k,w)=\bar c$.  Suppose for a moment the second Black term is
negligible; then $\PhiN(d_{+})=\bar c$, i.e.\
$d_{+}=q$ with $q=\PhiN^{-1}(\bar c)$ (a negative number --- on the
running node $\bar c=\MacWingEnvCbar$, far below one half).  Unpack
$d_{+}$:
\[
\frac{\sqrt w}{2}-\frac{k}{\sqrt w}=q
\qquad\Longrightarrow\qquad
w-2q\sqrt w-2k=0
\qquad\Longrightarrow\qquad
\sqrt{w^{+}(k)}=q+\sqrt{q^{2}+2k},
\]
taking the positive root of the quadratic in $\sqrt w$.  The neglected
term justifies itself with unexpected grace.  The two Black arguments
always differ by $\sqrt w$, so along this solution
$d_{-}=d_{+}-\sqrt w=q-\sqrt{w}$, and
\[
\frac{d_{-}^{2}}{2}
=\frac{q^{2}-2q\sqrt w+w}{2}
=\frac{q^{2}+2k}{2}
\qquad\Longrightarrow\qquad
\log\big(e^{k}\PhiN(d_{-})\big)
\le k-\frac{q^{2}+2k}{2}
=-\frac{k}{2}-\frac{q^{2}}{2},
\]
using the quadratic itself in the middle step, and the tail bound of
\cref{sec:vsceiling} (whose $-\log|d_{-}|$ term only strengthens the
inequality) in the last: the correction decays like $e^{-k/2}$,
exponentially small.  (At $k=30$ the closed form agrees with a full
Black inversion by the reference implementation --- the working code
behind every number in this book --- to \MacWingEnvTopAgree\ in
total variance.)  Differentiating the closed
form gives the edge's slope,
\begin{equation}
\frac{\dd w^{+}}{\dd k}
=2\big(q+\sqrt{q^{2}+2k}\big)\cdot\frac{1}{\sqrt{q^{2}+2k}}
=2+\frac{2q}{\sqrt{q^{2}+2k}}
\;\longrightarrow\;2
\qquad(k\to\infty),
\label{eq:vstopslope}
\end{equation}
rising monotonically (recall $q<0$, and the negative term shrinks) ---
it reads \MacWingEnvTopSlopeDeep\ at $k=30$ --- and approaching $2$
from below, never reaching it.  The top edge
of the finite-strike envelope is the ceiling of
\cref{prop:vsceiling}, met for the third time and in its sharpest
form: the steepest admissible completion climbs at slope
$2-2|q|/\sqrt{q^{2}+2k}$, a curve that grazes the ceiling
asymptotically while staying strictly under it --- the supremum, not
a member, made into a formula.

The section's summary is uncomfortable and important.  Between the
floor and the cap, mathematics is silent: every curve in the fan
prices a legitimate law consistent with everything quoted.  The data
cannot choose.  The model must --- and the next section is about
doing that in the open.
